\section{Mid-Term Examination --- Set A}
\label{sec:midterm_set_a}

\LPUPaperHeader{Mid-Term Examination (Objective MCQ)}{A}{90 Minutes}{30 Marks}{CO1 (Unit 1: Matrices), CO2 (Unit 2: ODEs \& Calculus), CO3 (Unit 3: Infinite Series)}

\begin{center}
\sffamily\textbf{\color{AlertRed}Instructions: All 30 questions are compulsory. Each question carries 1 mark. A penalty of 0.25 marks is deducted for each incorrect response.}
\end{center}

\vspace{3mm}

% ------------------------------------------------------------------------------
% 30 MCQS
% ------------------------------------------------------------------------------

\mcqitem{1}{The differential coefficient of $\log(\tan x)$ with respect to $x$ is:}{$2\sec 2x$}{$2\csc 2x$}{$2\sec^2 x$}{$2\csc^2 x$}{1}

\mcqitem{2}{The value of the indefinite integral $\int \frac{dx}{x^2 + 2x + 2}$ is equal to:}{$\tan^{-1}(x+1) + C$}{$\tan^{-1} x + C$}{$(x+1)\tan^{-1} x + C$}{$\frac{1}{2}\tan^{-1}(x+1) + C$}{1}

\mcqitem{3}{If $y = x^x$ for $x > 0$, then the derivative $\frac{dy}{dx}$ is:}{$x^x \ln x$}{$x^x(1 + \ln x)$}{$2 + \ln x$}{$x^{x-1}$}{1}

\mcqitem{4}{The derivative of $f(x) = \ln(x + \sqrt{x^2 + 1})$ with respect to $x$ is:}{$\sqrt{x^2 + 1}$}{$\frac{x}{\sqrt{x^2+1}}$}{$\frac{1}{\sqrt{x^2 + 1}}$}{$1 + \frac{x}{\sqrt{x^2+1}}$}{1}

\mcqitem{5}{The value of the definite integral $\int_0^\pi \sin^2 x \, dx$ is:}{$0$}{$\frac{\pi}{2}$}{$\pi$}{$\frac{\pi}{4}$}{1}

\mcqitem{6}{For a non-singular square matrix $A$ of order $3$ with $\det(A) = 3$, the value of $\det(\text{adj}(A))$ is:}{$3$}{$9$}{$27$}{$81$}{1}

\mcqitem{7}{If $A$ is an invertible square matrix, then $\det(A^{-1})$ is strictly equal to:}{$\frac{1}{\det(A)}$}{$\det(A)$}{$-\frac{1}{\det(A)}$}{$-\det(A)$}{1}

\mcqitem{8}{If $2, 3, -4$ are the eigenvalues of a matrix $A$, then the eigenvalues of the matrix $A^2$ are:}{$2, 3, -4$}{$-2, -3, 4$}{$4, 9, 16$}{$4, 3, 16$}{1}

\mcqitem{9}{The rank of the upper triangular matrix $A = \begin{bmatrix} 0 & 2 & -1 \\ 1 & 5 & 2 \\ 0 & 0 & 7 \end{bmatrix}$ is:}{$1$}{$2$}{$3$}{$0$}{1}

\mcqitem{10}{For any non-singular matrix $A$ of order $n$, the matrix $\text{adj}(\text{adj}(A))$ is equal to:}{$|A|^{n-1} A$}{$|A|^{n-2} A$}{$|A|^n A$}{$|A|^{n-1} I$}{1}

\mcqitem{11}{The value of $\int \frac{e^{\tan^{-1}x}}{1 + x^2} \, dx$ is equal to:}{$e^{\tan^{-1}x} + C$}{$\frac{(e^{\tan^{-1}x})^2}{2} + C$}{$\frac{(\tan^{-1}x)^2}{2} + C$}{$\tan^{-1}(e^x) + C$}{1}

\mcqitem{12}{The definite integral $\int_{-1}^1 (x^3 + x) \, dx$ evaluates to:}{$0$}{$\frac{\pi}{2}$}{$\pi$}{$1$}{1}

\mcqitem{13}{If $x^y = e^{x-y}$, then $\frac{dy}{dx}$ is equal to:}{$\frac{y}{(1+\ln x)^2}$}{$\frac{\ln x}{(1 + \ln x)^2}$}{$\frac{x}{(1+\ln x)^2}$}{$\frac{1}{(1+\ln x)^2}$}{1}

\mcqitem{14}{The eigenvalues of a Hermitian matrix are always:}{All Real}{All Purely Imaginary}{Zero}{Unitary}{1}

\mcqitem{15}{If $A$ is an $m \times n$ matrix where $m < n$, then the maximum possible value of $\text{rank}(A)$ is:}{$n$}{$m$}{$m \cdot n$}{$n - m$}{1}

\mcqitem{16}{The value of $\int \frac{e^x(1+x)}{\cos^2(x e^x)} \, dx$ is equal to:}{$-\cot(x e^x) + C$}{$\tan(x e^x) + C$}{$\sec(x e^x) + C$}{$\tan(e^x) + C$}{1}

\mcqitem{17}{The system of linear equations $AX = B$ is \textbf{inconsistent} (has no solution) if:}{$\text{rank}(A) = \text{rank}[A|B] = n$}{$\text{rank}(A) \neq \text{rank}[A|B]$}{$\text{rank}(A) = \text{rank}[A|B] < n$}{$\det(A) \neq 0$}{1}

\mcqitem{18}{The roots of the auxiliary equation for $(D^2 + 9)y = 0$ are:}{$\pm 3$}{$\pm 3i$}{$3, 3$}{$-3, -3$}{1}

\mcqitem{19}{The particular integral (P.I.) of $(D - 2)^2 y = e^{2x}$ is:}{$\frac{x^2}{2} e^{2x}$}{$x e^{2x}$}{$\frac{x}{2} e^{2x}$}{$\frac{x^3}{6} e^{2x}$}{1}

\mcqitem{20}{The general solution of the Cauchy-Euler equation $x^2 y'' - 2x y' + 2y = 0$ is:}{$c_1 x + c_2 x^2$}{$c_1 e^x + c_2 e^{2x}$}{$(c_1 + c_2 \ln x)x$}{$c_1 x + c_2 x^{-2}$}{1}

\mcqitem{21}{The infinite $p$-series $\sum_{n=1}^\infty \frac{1}{n^p}$ is \textbf{convergent} if and only if:}{$p \le 1$}{$p < 1$}{$p > 1$}{$p \ge 0$}{1}

\mcqitem{22}{Using D'Alembert's Ratio Test, the series $\sum_{n=1}^\infty u_n$ converges if $\lim_{n\to\infty} \left|\frac{u_{n+1}}{u_n}\right| = L$ satisfies:}{$L > 1$}{$L < 1$}{$L = 1$}{$L = \infty$}{1}

\mcqitem{23}{The alternating series $\sum_{n=1}^\infty (-1)^{n-1} \frac{1}{n} = 1 - \frac{1}{2} + \frac{1}{3} - \frac{1}{4} + \dots$ is:}{Divergent}{Absolutely Convergent}{Conditionally Convergent}{Oscillatory}{1}

\mcqitem{24}{The radius of convergence $R$ of the power series $\sum_{n=0}^\infty \frac{x^n}{n!}$ is:}{$0$}{$1$}{$e$}{$\infty$}{1}

\mcqitem{25}{The sum of the eigenvalues of the matrix $A = \begin{bmatrix} 4 & 2 & 1 \\ 0 & 3 & 5 \\ 0 & 0 & -2 \end{bmatrix}$ is:}{$5$}{$7$}{$-24$}{$10$}{1}

\mcqitem{26}{The product of the eigenvalues of the matrix $A = \begin{bmatrix} 2 & 1 \\ 4 & 5 \end{bmatrix}$ is:}{$6$}{$7$}{$10$}{$8$}{1}

\mcqitem{27}{The series $\sum_{n=1}^\infty \frac{1}{n^2 + 1}$ is:}{Convergent by Comparison with $\sum \frac{1}{n^2}$}{Divergent by Comparison with $\sum \frac{1}{n}$}{Oscillatory}{Conditionally Convergent}{1}

\mcqitem{28}{If $A$ is an orthogonal matrix, then $\det(A)$ is always:}{$0$}{$\pm 1$}{$\infty$}{Any real number}{1}

\mcqitem{29}{The Wronskian $W(y_1, y_2)$ of the functions $y_1 = \cos 2x$ and $y_2 = \sin 2x$ is:}{$1$}{$2$}{$-2$}{$0$}{1}

\mcqitem{30}{The particular integral of $(D^2 + 4)y = \cos 2x$ is:}{$-\frac{x}{4}\sin 2x$}{$\frac{x}{4}\sin 2x$}{$\frac{x}{2}\cos 2x$}{$-\frac{x}{2}\sin 2x$}{1}

\vspace{5mm}
\hrule
\vspace{5mm}

% ------------------------------------------------------------------------------
% ANSWER KEY & DETAILED EXPLANATIONS
% ------------------------------------------------------------------------------
\subsection*{Answer Key (Mid-Term Set A)}

\begin{center}
\begin{tabular}{|c|c||c|c||c|c||c|c||c|c||c|c|}
\hline
\textbf{Q} & \textbf{Ans} & \textbf{Q} & \textbf{Ans} & \textbf{Q} & \textbf{Ans} & \textbf{Q} & \textbf{Ans} & \textbf{Q} & \textbf{Ans} & \textbf{Q} & \textbf{Ans} \\
\hline
1 & \textbf{B} & 6 & \textbf{B} & 11 & \textbf{A} & 16 & \textbf{B} & 21 & \textbf{C} & 26 & \textbf{A} \\
2 & \textbf{A} & 7 & \textbf{A} & 12 & \textbf{A} & 17 & \textbf{B} & 22 & \textbf{B} & 27 & \textbf{A} \\
3 & \textbf{B} & 8 & \textbf{C} & 13 & \textbf{B} & 18 & \textbf{B} & 23 & \textbf{C} & 28 & \textbf{B} \\
4 & \textbf{C} & 9 & \textbf{C} & 14 & \textbf{A} & 19 & \textbf{A} & 24 & \textbf{D} & 29 & \textbf{B} \\
5 & \textbf{B} & 10 & \textbf{B} & 15 & \textbf{B} & 20 & \textbf{A} & 25 & \textbf{A} & 30 & \textbf{B} \\
\hline
\end{tabular}
\end{center}

\vspace{3mm}
\subsection*{Selected Detailed Mathematical Explanations}

\begin{solutionbox}[Explanations for Critical Questions]
\textbf{Q.1 (Ans: B)} $\frac{d}{dx}[\ln(\tan x)] = \frac{1}{\tan x}\sec^2 x = \frac{\cos x}{\sin x}\frac{1}{\cos^2 x} = \frac{1}{\sin x \cos x} = \frac{2}{2\sin x \cos x} = \frac{2}{\sin 2x} = 2\csc 2x$.\\[2mm]
\textbf{Q.6 (Ans: B)} For order $n=3$, $\det(\text{adj}(A)) = (\det A)^{n-1} = 3^{3-1} = 3^2 = 9$.\\[2mm]
\textbf{Q.8 (Ans: C)} By spectral mapping theorem, if $\lambda$ is an eigenvalue of $A$, then $\lambda^k$ is an eigenvalue of $A^k$. Hence $2^2=4, 3^2=9, (-4)^2=16$.\\[2mm]
\textbf{Q.10 (Ans: B)} Standard theorem: $\text{adj}(\text{adj}(A)) = |A|^{n-2} A$.\\[2mm]
\textbf{Q.13 (Ans: B)} $x^y = e^{x-y} \implies y \ln x = x - y \implies y(1 + \ln x) = x \implies y = \frac{x}{1 + \ln x}$.\\
Differentiating: $\frac{dy}{dx} = \frac{(1+\ln x)(1) - x(1/x)}{(1+\ln x)^2} = \frac{\ln x}{(1+\ln x)^2}$.\\[2mm]
\textbf{Q.19 (Ans: A)} $\frac{1}{(D-2)^2}e^{2x} = \frac{x^2}{2!}e^{2x} = \frac{x^2}{2}e^{2x}$.\\[2mm]
\textbf{Q.23 (Ans: C)} The alternating series converges by Leibniz test. However, the series of absolute values $\sum \frac{1}{n}$ is the harmonic series which diverges. Hence it is \textit{conditionally convergent}.\\[2mm]
\textbf{Q.24 (Ans: D)} Radius of convergence for $e^x = \sum \frac{x^n}{n!}$ is $R = \lim_{n\to\infty} \frac{(n+1)!}{n!} = \lim (n+1) = \infty$.\\[2mm]
\textbf{Q.25 (Ans: A)} Sum of eigenvalues = Trace of matrix $= 4 + 3 + (-2) = 5$.\\[2mm]
\textbf{Q.26 (Ans: A)} Product of eigenvalues = Determinant of matrix $= (2)(5) - (1)(4) = 10 - 4 = 6$.\\[2mm]
\textbf{Q.29 (Ans: B)} $W = \begin{vmatrix} \cos 2x & \sin 2x \\ -2\sin 2x & 2\cos 2x \end{vmatrix} = 2\cos^2 2x + 2\sin^2 2x = 2(1) = 2$.\\[2mm]
\textbf{Q.30 (Ans: B)} $\frac{1}{D^2+4}\cos 2x = \frac{x}{2(2)}\sin 2x = \frac{x}{4}\sin 2x$.
\end{solutionbox}

\newpage
